% Preamble additions needed for this chapter:
% \usepackage{listings}
% \lstset{basicstyle=\ttfamily\footnotesize, breaklines=true,
%         frame=single, columns=flexible, keepspaces=true,
%         showstringspaces=false, xleftmargin=0pt}

\chapter{Implementation}

This chapter describes how the testbed was built and how the experiments are run on
it. It covers the DTN and mail layers, the delay mechanism, the resolver, the setup
of each anti-spam tool, and the harness that executes a full set of runs without
supervision.

\section{Basis and Scope of This Work}

The DTN transport layer of the testbed follows the approach established in earlier work on delay-tolerant email, which built and validated an ION and BPMail pipeline for carrying mail over the Bundle Protocol. The same set of software is used, but instead of being containerised, the software stack is built on virtual machines, for the reasons given in the previous chapter.

All other components of the testbed architecture are new: delay introduction on the relay, the external resolver, installation and setup of all anti-spam tools, the measurement methodology, and the experimental design.

\section{Environment and Dependencies}

The testbed uses five VMs on an Apple Silicon MacBook Air, under UTM. All five run Ubuntu 22.04 LTS (ARM64). Three servers are named mail1, space and mail2; two clients, one for each network. The mail, DTN and anti-spam applications are installed on the servers. Clients run Thunderbird and serve only to verify that mail transfer works end-to-end; they do not participate in measurement.

Three libraries are more recent than those in the Ubuntu 22.04 repositories and must be built from source on each server before ION and BPMail are compiled: mbedtls 2.28.8, gmime 3.2.7 and c-ares 1.34.4. Each is built in order, then \texttt{ldconfig} is run so the libraries are found first.

\section{DTN Layer}

ION-DTN 4.1.4 is built from source on all three servers:

\begin{lstlisting}
./configure --disable-crypto CFLAGS="-Dlinux"
make
sudo make install
sudo ldconfig
\end{lstlisting}

The \texttt{-Dlinux} flag is required because ION's \texttt{platform.h} selects the
system headers that define \texttt{MAXPATHLEN} behind an \texttt{\#ifdef linux}
check. Crypto is disabled because the Bundle Protocol Security extensions are not
used here.

Each node is configured from the \texttt{bench-udp} demo template supplied with ION.
The node numbers are 1 for mail1, 2 for space and 3 for mail2. Bundles are carried
over the UDP convergence layer, with each node listening on its own port: 1113 on
mail1, 2113 on space and 3113 on mail2. The bundle protocol configuration for mail1
is:

\begin{lstlisting}
a scheme ipn 'ipnfw' 'ipnadminep'
a endpoint ipn:1.0 x
a endpoint ipn:1.1 x
a endpoint ipn:1.128 x
a endpoint ipn:1.129 x
a protocol udp 1400 100
a induct  udp 192.168.100.10:1113 udpcli
a outduct udp 192.168.100.1:2113  udpclo
\end{lstlisting}

Space is the only node with two inducts and two outducts, and is the only route
between the two networks. Routing uses Contact Graph Routing, so routes are computed
from the contact plan rather than configured statically. The contact plan is
identical on all three nodes and declares reachability in both directions between
each adjacent pair:

\begin{lstlisting}
m horizon +0
a range   +0 +86400  1 2  1
a range   +0 +86400  2 1  1
a range   +0 +86400  2 3  1
a range   +0 +86400  3 2  1
a contact +0 +86400  1 2  100000000
a contact +0 +86400  2 1  100000000
a contact +0 +86400  2 3  100000000
a contact +0 +86400  3 2  100000000
\end{lstlisting}

The contact windows are set to 24 hours. Shorter windows expire during a run and
bundles stop being forwarded, so the window has to cover a full pass of three runs
with margin. Each node is started with an \texttt{ionstart} script that loads the
node configuration, the contact plan, the security database and the bundle protocol
configuration in that order, with short pauses between them.

\section{Mail Layer}

BPMail 0.1.0 provides the gateway between the mail system and ION. It is installed
on mail1 and mail2 only; space forwards bundles and does not handle mail.

\subsection{Build}

BPMail is built with Meson. One change to the build configuration is required. The
first build failed with an error reporting that no value was defined for
\texttt{MAXPATHLEN}. The cause is an interaction between two settings. ION's
\texttt{platform.h} includes the system headers that define \texttt{MAXPATHLEN} only
when the \texttt{linux} platform macro is set, guarded by \texttt{\#ifdef linux}.
BPMail's \texttt{meson.build} specified \texttt{c\_std=c17}, and strict ISO C17
disables non-standard platform macros such as \texttt{linux}, so the guard did not pass and the definition was never picked up, even when \texttt{-Dlinux} was passed
on the command line. Changing the standard to GNU C17 keeps those platform macros
while remaining C17-compliant:

\begin{lstlisting}
sed -i "s/c_std=c17/c_std=gnu17/" ~/bpmail/meson.build
\end{lstlisting}

The platform flag is supplied through a native file, and the build then completes:

\begin{lstlisting}
printf '[built-in options]\nc_args = [\x27-Dlinux\x27]\n' > linux-native.txt
meson setup build --native-file linux-native.txt
ninja -C build
sudo ninja -C build install
\end{lstlisting}

\subsection{DTPC configuration}

BPMail transfers mail over DTPC rather than raw bundles, so the two mail nodes need
DTPC endpoints and a profile. Service 128 is used for sending and 129 for receiving,
and the endpoints are added to each node's bundle protocol configuration as shown
above. The DTPC profile is the same on both nodes:

\begin{lstlisting}
1
w 1
a profile 1 5 1000 10 600 0.1 dtn:none
s
\end{lstlisting}

The profile sets a maximum of five retransmissions, an aggregation size limit of
1000 bytes, an aggregation time limit of 10 seconds and a bundle lifetime of 600
seconds. The \texttt{dtpcadmin} step is appended to the \texttt{ionstart} script on
both mail nodes so that DTPC comes up with the rest of ION.

A DTPC topic can only be opened once per node, for sending or receiving but not
both, so each direction uses its own topic. The experiments send from mail1 to
mail2, and use topic 26 for that direction. The destination endpoint is the
receiving service on node 3, \texttt{ipn:3.129}; sending to the service-128 endpoint
does not reach the receiver.

\subsection{Postfix and Dovecot}

Postfix is configured on both mail nodes with a transport map that routes mail for
the remote domain to a pipe transport, which hands the message to BPMail. BPMail
needs access to ION's shared memory, which depends on the working directory, so the
pipe calls a small wrapper that changes to the ION node directory first:

\begin{lstlisting}
bpmail   unix  -  n  n  -  1  pipe
  flags=F user=<user> argv=/home/<user>/bpmailsend_wrapper.sh 1 ipn:3.129
\end{lstlisting}

A background daemon runs \texttt{bpmailrecv} in a loop on each node and pipes what it
receives into \texttt{sendmail}, which delivers it to the local mailbox. Dovecot
serves those mailboxes over IMAP, in mbox format.

This path is what the Thunderbird clients use, and how the testbed was verified as a working mail system. It is not the path used for measurement. The experiments call bpmailsend and bpmailrecv directly, for the reasons given in Section~\ref{sec:harness-agents}.

\section{Delay Simulation}

The link delay is applied on space with tc and the netem queueing discipline, on the Moon-facing interface, so it acts on all traffic crossing the boundary between the two networks. As explained in \ref{sec:delay}, tc shapes only egress, so the outbound leg uses a netem qdisc on the interface root and the inbound leg is redirected to an ifb device carrying the same delay. The full sequence is:

\begin{lstlisting}
sudo modprobe ifb
sudo ip link add ifb0 type ifb
sudo ip link set ifb0 up
sudo tc qdisc add dev enp0s2 root netem delay 1300ms
sudo tc qdisc add dev enp0s2 handle ffff: ingress
sudo tc filter add dev enp0s2 parent ffff: protocol ip u32 match u32 0 0 \
     action mirred egress redirect dev ifb0
sudo tc qdisc add dev ifb0 root netem delay 1300ms
\end{lstlisting}

Clearing the delay removes both qdiscs, the ingress filter and the ifb device.

The one-way delay is a parameter. The experiments use 1300\,ms in each direction,
giving a round-trip time of about 2.6\,s. Because the delay sits at the boundary of
the Moon network, everything mail2 exchanges with anything outside it pays that
cost: the bundles it receives from mail1, and the DNS queries the scanner makes
while examining a message.

\section{DNS Resolver}

The scanners run on mail2 and need working DNS. The queries are served by a local
\texttt{unbound} instance on mail2, which forwards to a private recursive resolver
running on an Azure virtual machine on port 5353.

Two changes to the default \texttt{unbound} configuration were needed.

The first was a startup failure. With DNSSEC validation enabled and all queries
forwarded, \texttt{unbound} logged repeated failures to prime the trust anchor and
refused to resolve anything. Forwarding and validation conflict in this
configuration. Because the upstream resolver performs validation, the validator
module can be removed on mail2, leaving only the iterator:

\begin{lstlisting}
module-config: "iterator"
\end{lstlisting}

The second concerns two timers that interact with the link delay:

\begin{lstlisting}
discard-timeout: 9000
infra-cache-min-rtt: 4000
\end{lstlisting}

Both defaults are shorter than the round-trip time of the delayed link. The reason
these values are needed, and what happens with the defaults in place, is a result of
the experiments and is presented in Chapter~5.

\section{Anti-Spam Tools}

Three tools are installed on mail2, the node that receives and scans. Each is
invoked once per message, on the message file, and each is run under two
configurations: a default one, and one with its timeout reduced below the round-trip
time of the link.

\subsection{SpamAssassin}

SpamAssassin is invoked through its standalone binary rather than through
\texttt{spamc}:

\begin{lstlisting}
spamassassin < "$FILE"
\end{lstlisting}

The client and daemon pair was tried first and had to be abandoned.When
\texttt{spamc} cannot reach \texttt{spamd} it fails open, passing the message through
unchanged, so every message came back with no \texttt{X-Spam-Status} header and no
indication that anything had gone wrong. The standalone binary scans without a
daemon, and re-reads \texttt{local.cf} on every invocation, which means a
configuration change takes effect on the next message with no restart.

The two configurations differ only in whether the timeouts are set. The default
profile removes any timeout lines from \texttt{local.cf}, leaving SpamAssassin's own
defaults of 15\,s for \texttt{rbl\_timeout} and 5\,s for \texttt{spf\_timeout}. The
reduced profile sets both to 1.5\,s, below the 2.6\,s round-trip time:

\begin{lstlisting}
sudo sed -i '/^rbl_timeout /d;/^spf_timeout /d' /etc/spamassassin/local.cf
echo 'rbl_timeout 1.5' | sudo tee -a /etc/spamassassin/local.cf
echo 'spf_timeout 1.5' | sudo tee -a /etc/spamassassin/local.cf
\end{lstlisting}

\subsection{rspamd}

rspamd runs as a daemon and is invoked through its client:

\begin{lstlisting}
rspamc --ip 1.2.3.4 --mime < "$FILE"
\end{lstlisting}

The \texttt{--ip} option supplies a client address so the message is treated as having arrived from outside. Without it, mail presented locally is treated as internal and the network checks central to this study are skipped.

rspamd caches state in the running daemon, so it is restarted before each run to
start from a cold cache. The reduced-timeout profile writes a single override file
and restarts:

\begin{lstlisting}
echo 'task_timeout = 1.5s;' | sudo tee /etc/rspamd/local.d/options_override.inc
sudo systemctl restart rspamd
\end{lstlisting}

The default profile removes that file and restarts.

\subsection{Vipul's Razor}

Razor is invoked with \texttt{razor-check}, which reports its verdict through the
exit code: 0 if the message is catalogued as spam, 1 if it is not, and 2 on a network
error. The third case matters here, because it separates a decision from a failure to
reach the server:

\begin{lstlisting}
razor-check < "$FILE"
\end{lstlisting}

Razor's server list is refreshed with \texttt{razor-admin -discover} before each run.

Razor differs from the other two in one respect that affects the experiment
directly: it has no timeout setting in \texttt{razor-agent.conf}. Its connection
timeouts are fixed in the client source. Producing the reduced-timeout condition
therefore requires patching the module that holds them, keeping a copy of the
original so the default condition can be restored:

\begin{lstlisting}
sudo cp -n "$CORE_PM" "$CORE_PM.orig"
sudo sed -i 's/Timeout  *=>  *20/Timeout => 2/g' "$CORE_PM"
\end{lstlisting}

That a source patch is the only way to change this value is itself a finding about
the tool, and is returned to in Chapter~5.

\section{Email Corpus}
The corpus is drawn from Richard Clayton's public spam archive, a collection of spam messages captured from live mail traffic. 681 of these were used in every run. The same 681 messages, in the same sorted order, were scanned in all three runs of every tool, so per-message comparisons across runs refer to identical inputs. The corpus is used to characterise tool behaviour under delay rather than to measure absolute detection accuracy, so what matters is that the message set is fixed and realistic, not that it is balanced between spam and legitimate mail.

\section{Run Harness}
\label{sec:harness}

The three runs for one tool take several hours and involve three machines, so the
experiments are automated. The harness is a Python program that runs on the Mac and
drives the testbed over SSH. A single command performs a complete pass:

\begin{lstlisting}
python3 run.py --tool spamassassin
\end{lstlisting}

Every parameter it uses is held in one configuration file, so the procedure is identical between tools and between passes, and the file itself records the exact settings under which a set of results was produced.

\subsection{Structure}

The harness connects to the three servers with Paramiko. Mail1 and mail2 are not
reachable directly from the Mac, so their connections are tunnelled through space,
which is opened first.

Work on the servers is done by two shell agents, one on mail1 that sends and one on
mail2 that receives and scans. They are written out to each node at the start of a
run, together with an environment file holding that run's settings, and launched
detached so that they survive the SSH session. Each agent reports through two files
in its working directory: \texttt{STATUS}, which holds \texttt{RUNNING},
\texttt{DONE} or \texttt{FAILED}, and \texttt{progress}, which holds a count of
messages handled. The host polls these every thirty seconds and prints a line for
each. The alternative, holding a live channel open to each node for the length of a
run, is fragile over hours; sentinel files survive a dropped connection.

\subsection{Preflight and boot}

Before anything starts, the harness checks that \texttt{utmctl} is usable, that the
named virtual machines exist, and that the results directory has at least 20\,GB
free. It then starts the machines in order and waits for each to accept SSH, up to
five minutes. It also holds \texttt{caffeinate} open for the session so the Mac does
not sleep partway through a run.

\subsection{Verification gates}

A run only starts if the testbed is in a known state. Four checks are made, and any
of them can abort the pass:

\begin{itemize}
  \item \textbf{Clocks.} Timing is compared between the sender on mail1 and the
  receiver on mail2, so their clocks must agree. The offset from the host is measured
  on each node; anything above one second aborts.
  \item \textbf{Resolver.} \texttt{unbound} must be active on mail2 and must answer a
  test query.
  \item \textbf{ION.} ION is restarted across all three nodes, which resets the
  contact window so that it covers the whole pass.
  \item \textbf{Tool.} The tool's version command must return output, confirming it is installed; the version string is recorded with the results.
\end{itemize}

\subsection{Sequence of a run}
\label{sec:harness-agents}

Each run follows the same sequence.

The testbed is first returned to a cold state: \texttt{unbound} is restarted, which
clears both its cache and its record of upstream server behaviour, any agent state
from a previous run is deleted, the tool is reset, and any leftover \texttt{netem}
configuration is removed. The resolver's infrastructure cache is dumped and saved, so
that its state before the run is on record.

The tool profile for the run is then applied, along with the delay if the run calls for it. The round-trip time is measured by pinging mail1 from mail2 and checked against a window: 2400 to 2900 ms for a delayed run, and under 150 ms for the baseline. A measurement outside the window aborts the run. This is the most important gate in the harness. A run whose delay was not actually in place, or was applied twice, would produce plausible-looking numbers that mean nothing, and the failure would not be obvious afterwards.

Packet capture is then started on space and on mail2, excluding management SSH so
that the harness's own traffic does not appear in the bandwidth figures:

\begin{lstlisting}
sudo tcpdump -i any -s 0 -w {out} -C 100 -W 30 'not tcp port 22'
\end{lstlisting}

The receiver is launched first and the sender three seconds later, so that no message
can arrive before the receiver is listening.

The sender walks the corpus in sorted order. For each message it records the
Message-ID and a timestamp, sends it, and records the exit code:

\begin{lstlisting}
mid=$(grep -i -m1 '^Message-ID:' "$f" | tr -d '\r\n' | cut -d' ' -f2-)
sent=$(date -u +%s.%N)
eval "$SEND_CMD"
rc=$?
echo "$seq,$bn,$mid,$sent,$rc" >> "$SENT"
\end{lstlisting}

The send command pipes the message directly to BPMail:

\begin{lstlisting}
sed "1{/^From /d}" "$FILE" | bpmailsend -t 26 1 ipn:3.129
\end{lstlisting}

The leading \texttt{From } line of the mbox-format corpus file is stripped, because
it is not part of the message.

The receiver blocks on \texttt{bpmailrecv} until a bundle arrives, writes the message
out, scans it, and records the timing:

\begin{lstlisting}
RAW=$(timeout "$RECV_TIMEOUT" bash -c "$RECV_CMD")
arrived=$(date -u +%s.%N)
printf '%s' "$RAW" > "$eml"
start=$(date -u +%s.%N)
eval "$SCAN_CMD" > "$outf" 2>&1
rc=$?
end=$(date -u +%s.%N)
\end{lstlisting}

The scan output is saved in full for every message, not just the parsed verdict, so
that the parsing can be revised later without repeating a run.

Sending and receiving through BPMail directly, rather than through Postfix, gives a
clean boundary for the scan time. The clock starts when the message arrives and stops
when the tool exits, so the recorded time is the time the tool spent on that message.
Invoking the tool through Postfix would include its queueing in every measurement,
and at a round-trip time of 2.6\,s the two could not be separated.

Messages are paced rather than sent back to back. The interval differs by run and is
set in the configuration.

\subsection{Completion and collection}

When the sender reports \texttt{DONE} the harness enters a drain phase, because
messages are still in flight and will keep arriving for at least one one-way delay
afterwards. It waits for the receiver's count to reach the number sent, up to fifteen
minutes, then writes a \texttt{STOP} file that the receiver checks between bundles, so it exits cleanly rather than being killed mid-scan.

Captures are stopped, and the resolver's state is recorded again. Two health figures
are taken from it: the highest retransmit timeout in the infrastructure cache, and
the number of log entries reporting answers discarded for arriving too late. A run
whose highest retransmit timeout has reached 120{,}000\,ms is marked
\texttt{SUSPECT} in its manifest, because that is the ceiling and indicates the
resolver stopped working properly during the run. The check runs automatically, so a compromised run is flagged at the time rather than discovered later.

All artifacts are then pulled back to the Mac. The sent and received records are
reconciled by Message-ID to identify any messages that did not arrive, and the
receiver's timing records are joined with the parsed scan output into a single
per-message file. Each tool has its own parser, because the three report their
verdicts differently: SpamAssassin in an \texttt{X-Spam-Status} header that folds
across lines, rspamd in \texttt{X-Spam-*} headers added to the message, and Razor in
its exit code alone.

Each run finally writes a manifest recording the tool and version, the profile, the measured round-trip time and the gate against which it was checked, the timings, the message counts, the resolver health figures, and the resulting status. The manifest is what makes a set of results interpretable months later.

\section{Running an Experiment}

A pass consists of three runs against the same corpus of 681 messages:

\begin{itemize}
  \item \textbf{Run 1}, no delay, default tool configuration. The baseline.
  \item \textbf{Run 2}, delay applied, default tool configuration.
  \item \textbf{Run 3}, delay applied, tool timeout reduced below the round-trip
  time.
\end{itemize}

A completed pass leaves, for each run, the sender's record, the receiver's timing
record, every scan output, packet captures from space and mail2, the resolver state
before and after, the joined per-message file, and the manifest. A summary across the
three runs is written at the end of the pass.

\section{Summary}

The testbed is a working DTN mail system with a controllable link delay, a resolver that reaches live blocklists, and three anti-spam tools, each configurable into a default and a reduced-timeout condition. The experiments are driven by a configuration-file-based harness that verifies the testbed before each run, refuses to proceed if the delay is not correct, collects every artifact, and flags runs whose resolver state indicates the results cannot be trusted.